\documentclass{article} % say
\usepackage{amsmath}
\usepackage{gensymb}
\usepackage{hyperref}
\usepackage{tikz}
\usepackage{tkz-euclide}
\usetikzlibrary {babel,angles,quotes,positioning}
\usepackage[utf8]{inputenc}
\usepackage[russian]{babel}
\begin{document}

\textbf{Задача}: Дано два отрезка: $a$ и $d$. С помощью циркуля и линейки постройте прямугольный треугольник, в котором длина одного из катетов равна длине $a$ и разница длин гипотезы и второго катета равна длине $d$.

\textbf{Решение}: Обозначим длину второго катета за $x$, тогда:

\begin{align*}
  a^2 + x^2 &= (x+d)^2 \\
  a^2 + x^2 &= x^2 + 2xd + d^2 \\
  a^2 - d^2 &= 2xd \\
  (a - d) (a + d) &= 2xd
\end{align*}

Таким образом:

\begin{equation}
  x = \frac{a + d}{2d} \cdot (a - d) 
\end{equation}

Построим треугольник $ABC$, причём $|AB| = a + d$ и $|BC| = 2d$, с любым углом между ними, например с прямым.
На отрезке $BC$ найдём точку $D$, такую что $|BD| = a - d$.
Через точку $D$ проведём прямую, параллельную $AC$. Точку пересечения этой прямой с $BC$ обозначим за $E$.

\begin{center}
\begin{tikzpicture}[scale=1.25]
\tkzDefPoints{0/6/A,0/0/B,4/0/C,2/0/D}
\tkzDefLine[parallel=through D](A,C) \tkzGetPoint{d'}
\tkzInterLL(B,A)(D,d') \tkzGetPoint{E}
\tkzDrawLine[add=0.2 and 0.2](D,E)

\tkzDrawPolygon(A,B,C)

\tkzMarkAngles[mark=|,size=0.5](B,A,C B,E,D)
\tkzMarkAngles[mark=||,size=0.5](E,D,B A,C,D)

\tkzDrawPoints(A,B,C,D,E)
\tkzLabelPoints(B,C)
\tkzLabelPoints[right](A,E)
\tkzLabelPoints[above right](D)
\end{tikzpicture}
\end{center}

Треугольник $BED$ по построение подобен треугольник $BAC$. А значит $\frac{|BE|}{|BD|} = \frac{|BA|}{|BC|}$, из чего:

\begin{equation}
  |BE| = \frac{|BA|}{|BC|} \cdot |BD| = \frac{a + d}{2d} \cdot (a - d) = x
\end{equation}

Теперь осталось построить прямоугольный треугольник по двум катетам.

\textbf{Источник}: А.~П.~Киселёв, <<Элементарная геометрия. Книга для учителя>> (1980), страница 36, задача 44. \url{https://disk.yandex.ru/d/QqFzGPzxTVXGag}

\end{document}


